\begin{tcolorbox}[breakable,colback=blue!5!white, colframe=blue!75!black, title=Seed Idea - \texttt{batch\_size\_grokking}]
\small
\begin{verbatim}
"Name": "batch_size_grokking",
"Title": "Batch Size Grokking: Assessing the impact of the training
batchsize on the grokking phenomenon",
"Experiment": "Modify the experiments to dynamically adjust the batch size
during training, starting with a small batch size and gradually increasing
it. This could potentially lead to faster generalization on the validation
set.",
"Interestingness": 6,
"Feasibility": 4,
"Novelty": 4,
"novel": true
\end{verbatim}
\end{tcolorbox}

\begin{tcolorbox}[breakable,colback=blue!5!white, colframe=blue!75!black, title=Idea 1/50 - \texttt{model\_size\_grokking}]
\small
\begin{verbatim}
"Name": "model_size_grokking",
"Title": "Investigating the Impact of Model Size on the Grokking
Phenomenon",
"Experiment": "Modify the Transformer class to accept variable number of
layers and dimension sizes. Test models with 1, 2, 4, and 8 layers, and
dimension sizes of 64, 128, 256, and 512. For each dataset and model size,
track the step at which grokking occurs (defined as validation accuracy
exceeding 99%
and final accuracy vs model size for each task.",
"Interestingness": 8,
"Feasibility": 7,
"Novelty": 7,
"novel": true
\end{verbatim}
\end{tcolorbox}

\begin{tcolorbox}[breakable,colback=blue!5!white, colframe=blue!75!black, title=Idea 2/50 - \texttt{optimizer\_grokking}]
\small
\begin{verbatim}
"Name": "optimizer_grokking",
"Title": "Optimization Dynamics and Grokking: Comparing SGD and Adam with
Different Learning Rate Schedules",
"Experiment": "Modify the training loop to support two optimizers (SGD,
Adam) and two learning rate schedules (constant, cosine annealing). For
each combination, run multiple experiments with different random seeds.
Track validation accuracy, training loss, and L2 norm of weight updates
throughout training. Compare the timing and extent of grokking across these
optimization strategies for each dataset. Analyze how different
optimization dynamics correlate with grokking behavior, including
statistical analysis of the results.",
"Interestingness": 9,
"Feasibility": 8,
"Novelty": 8,
"novel": true
\end{verbatim}
\end{tcolorbox}

\begin{tcolorbox}[breakable,colback=blue!5!white, colframe=blue!75!black, title=Idea 3/50 - \texttt{biased\_data\_grokking}]
\small
\begin{verbatim}
"Name": "biased_data_grokking",
"Title": "Grokking Under Biased Data: The Effect of Input Range Bias on
Neural Network Generalization",
"Experiment": "Modify the fetch_train_example method in AbstractDataset to
introduce a simple bias: favoring lower-valued inputs. For modular
arithmetic operations, sample 70%
of the input range. For permutations, favor permutations with more elements
in their original positions. Keep the validation set unbiased. Run
experiments comparing grokking behavior on biased vs. unbiased training
sets. Track metrics such as steps to 99%
validation accuracy, and training loss. Analyze how this bias affects
grokking across different operations.",
"Interestingness": 8,
"Feasibility": 8,
"Novelty": 8,
"novel": true
\end{verbatim}
\end{tcolorbox}

\begin{tcolorbox}[breakable,colback=blue!5!white, colframe=blue!75!black, title=Idea 4/50 - \texttt{adaptive\_noise\_grokking}]
\small
\begin{verbatim}
"Name": "adaptive_noise_grokking",
"Title": "Adaptive Noise in Grokking: Investigating Input Perturbations on
Algorithmic Learning and Representations",
"Experiment": "Modify the GroupDataset class to add operation-specific
noise during training: (1) For modular arithmetic, add small integer
perturbations. (2) For permutations, occasionally swap two elements.
Implement three noise levels (low, medium, high) for each operation.
Compare grokking behavior across noise levels and operations, tracking
steps to 99%
loss. Analyze learned representations by visualizing attention patterns and
performing principal component analysis (PCA) on hidden states at different
training stages. Compare these representations between noisy and non-noisy
training to understand how noise affects the abstraction of concepts.",
"Interestingness": 9,
"Feasibility": 8,
"Novelty": 9,
"novel": true
\end{verbatim}
\end{tcolorbox}

\begin{tcolorbox}[breakable,colback=blue!5!white, colframe=blue!75!black, title=Idea 5/50 - \texttt{attention\_evolution\_grokking}]
\small
\begin{verbatim}
"Name": "attention_evolution_grokking",
"Title": "Attention Evolution in Grokking: Quantifying the Transition from
Memorization to Generalization",
"Experiment": "Modify the Transformer class to output attention weights.
Extract and store attention weights at key checkpoints: start, mid-
training, grokking point (99%
visualization tools for attention heatmaps and create plots showing
attention evolution over time. Calculate the Frobenius norm of the
difference between attention matrices at consecutive checkpoints to
quantify attention evolution. Compare attention patterns and evolution
metrics across different operations (modular arithmetic vs. permutations).
Analyze attention for specific, informative input sequences to enhance
interpretability. Correlate attention evolution metrics with validation
accuracy and generalization performance.",
"Interestingness": 9,
"Feasibility": 8,
"Novelty": 8,
"novel": true
\end{verbatim}
\end{tcolorbox}

\begin{tcolorbox}[breakable,colback=blue!5!white, colframe=blue!75!black, title=Idea 6/50 - \texttt{local\_vs\_global\_attention\_grokking}]
\small
\begin{verbatim}
"Name": "local_vs_global_attention_grokking",
"Title": "Local vs Global Attention: Investigating the Impact of Attention
Scope on Grokking in Algorithmic Learning",
"Experiment": "Modify the DecoderBlock class to support two attention
mechanisms: full (global) attention and local attention with a fixed window
size. Implement these variants and run experiments across all datasets.
Track metrics including time to grokking (99%
validation accuracy, and training loss. Calculate and compare 'attention
entropy' for both mechanisms across tasks to quantify attention focus.
Analyze how the scope of attention (local vs global) affects grokking
behavior and final performance for different algorithmic tasks.",
"Interestingness": 9,
"Feasibility": 8,
"Novelty": 8,
"novel": true
\end{verbatim}
\end{tcolorbox}

\begin{tcolorbox}[breakable,colback=blue!5!white, colframe=blue!75!black, title=Idea 7/50 - \texttt{input\_encoding\_grokking}]
\small
\begin{verbatim}
"Name": "input_encoding_grokking",
"Title": "Binary vs One-Hot Encoding: Impact on Grokking in Algorithmic
Learning Tasks",
"Experiment": "Modify the AbstractDataset class to support two encoding
schemes: one-hot (current) and binary. Implement binary encoding for
modular arithmetic operations using log2(p) bits, and for permutations
using ceil(log2(k!)) bits to represent each permutation uniquely. Adjust
the Transformer class to accommodate different input sizes. Run experiments
for each encoding scheme across all datasets, tracking metrics such as time
to grokking (99%
loss, and model memory usage. Analyze how different encoding schemes affect
grokking behavior, convergence speed, and final performance for various
algorithmic tasks. Compare the impact of input representations on the
model's ability to learn and generalize across different operations.
Discuss how findings could inform input representation choices in other
machine learning tasks beyond algorithmic learning.",
"Interestingness": 9,
"Feasibility": 8,
"Novelty": 8,
"novel": true
\end{verbatim}
\end{tcolorbox}

\begin{tcolorbox}[breakable,colback=blue!5!white, colframe=blue!75!black, title=Idea 8/50 - \texttt{curriculum\_learning\_grokking}]
\small
\begin{verbatim}
"Name": "curriculum_learning_grokking",
"Title": "Curriculum Learning in Grokking: The Effect of Structured Example
Progression on Algorithmic Learning",
"Experiment": "Modify the AbstractDataset class to implement a simple
curriculum learning strategy. For modular arithmetic operations, start with
operations involving numbers in the lower half of the range and gradually
introduce larger numbers. For permutations, begin with permutations that
differ from the identity by one swap and progressively increase the number
of swaps. Implement a curriculum scheduler that increases difficulty every
500 steps. Run experiments comparing standard random sampling vs.
curriculum learning across all datasets. Track metrics including time to
grokking (99%
loss. Plot learning trajectories (validation accuracy over time) for both
approaches. Compare attention patterns between curriculum and random
approaches at different stages of training. Analyze how the curriculum
affects the grokking phenomenon across different operations and discuss
implications for training neural networks on algorithmic tasks.",
"Interestingness": 9,
"Feasibility": 8,
"Novelty": 8,
"novel": true
\end{verbatim}
\end{tcolorbox}

\begin{tcolorbox}[breakable,colback=blue!5!white, colframe=blue!75!black, title=Idea 9/50 - \texttt{weight\_init\_grokking}]
\small
\begin{verbatim}
"Name": "weight_init_grokking",
"Title": "Weight Initialization Strategies and Their Impact on Grokking in
Algorithmic Learning",
"Experiment": "Modify the Transformer class to support three weight
initialization strategies: Xavier/Glorot, Kaiming/He, and random normal (as
baseline). Implement these initialization methods for linear layers and
embeddings. Run experiments across all datasets for each initialization
strategy. Track metrics including time to grokking (99%
accuracy), final validation accuracy, training loss, and gradient norm
during training. Plot learning curves and compare the distribution of
weight values at different stages of training. Analyze the loss landscape
by computing gradient variance as a proxy for local geometry at key points
during training. Compare how different initialization strategies affect the
grokking phenomenon, convergence speed, and final performance across
various algorithmic tasks. Investigate potential correlations between
initial weight distributions, gradient variance characteristics, and the
timing/nature of the grokking transition.",
"Interestingness": 9,
"Feasibility": 9,
"Novelty": 8,
"novel": true
\end{verbatim}
\end{tcolorbox}

\begin{tcolorbox}[breakable,colback=blue!5!white, colframe=blue!75!black, title=Idea 10/50 - \texttt{task\_complexity\_grokking}]
\small
\begin{verbatim}
"Name": "task_complexity_grokking",
"Title": "Grokking Across Task Complexity: Mapping Neural Network Learning
Dynamics to Algorithmic Difficulty",
"Experiment": "1. Modify the AbstractDataset class to include new
operations of increasing complexity: modular addition, subtraction,
multiplication, and exponentiation. 2. Implement these operations in new
dataset classes. 3. Quantify task complexity using metrics like algebraic
degree and average solution time for humans (estimated). 4. Run experiments
for each operation, tracking metrics such as time to grokking (99%
validation accuracy), final validation accuracy, training loss, and a new
'complexity-adjusted learning rate' (validation accuracy improvement per
epoch, normalized by task complexity). 5. Plot learning curves and
complexity-adjusted learning rates for each operation. 6. Analyze attention
patterns and hidden state representations at different stages of training
for each operation. 7. Investigate correlations between quantified task
complexity and grokking characteristics (e.g., time to grokking, steepness
of accuracy improvement).",
"Interestingness": 9,
"Feasibility": 8,
"Novelty": 8,
"novel": true
\end{verbatim}
\end{tcolorbox}

\begin{tcolorbox}[breakable,colback=blue!5!white, colframe=blue!75!black, title=Idea 11/50 - \texttt{regularization\_grokking}]
\small
\begin{verbatim}
"Name": "regularization_grokking",
"Title": "The Role of Regularization in Grokking: How L2 and Label
Smoothing Affect Algorithmic Learning",
"Experiment": "1. Implement L2 regularization by adding weight decay to the
optimizer. 2. Implement label smoothing in the loss function. 3. Modify the
training function to support these regularization techniques with two
strength levels each (low and high). 4. Run experiments for each
regularization technique and strength across all datasets, including a
baseline without regularization. 5. Track metrics: time to grokking (99%
validation accuracy), final validation accuracy, training loss, and a new
'grokking speed' metric (rate of validation accuracy improvement from 50%
to 90%
for different regularization settings. 7. Analyze how L2 regularization and
label smoothing affect the timing, speed, and nature of grokking across
various algorithmic tasks, comparing against the non-regularized
baseline.",
"Interestingness": 9,
"Feasibility": 9,
"Novelty": 8,
"novel": true
\end{verbatim}
\end{tcolorbox}

\begin{tcolorbox}[breakable,colback=blue!5!white, colframe=blue!75!black, title=Idea 12/50 - \texttt{grokking\_extrapolation}]
\small
\begin{verbatim}
"Name": "grokking_extrapolation",
"Title": "Grokking and Extrapolation: Investigating the Limits of
Algorithmic Understanding",
"Experiment": "1. Modify AbstractDataset to create a separate test set with
out-of-distribution examples (e.g., larger numbers for modular arithmetic,
longer permutations). 2. Implement a new evaluation function for the test
set. 3. During training, periodically evaluate on both validation and test
sets. 4. Track metrics: time to grokking on validation set, final
validation accuracy, test set accuracy at grokking point, final test set
accuracy, and 'extrapolation gap'. 5. Implement visualization of test set
performance and extrapolation gap over time, highlighting the grokking
point. 6. Compare extrapolation capabilities across different operations
and model sizes. 7. Analyze attention patterns on test set inputs before
and after grokking. 8. Implement a simple MLP baseline for comparison.",
"Interestingness": 9,
"Feasibility": 8,
"Novelty": 9,
"novel": true
\end{verbatim}
\end{tcolorbox}

\begin{tcolorbox}[breakable,colback=blue!5!white, colframe=blue!75!black, title=Idea 13/50 - \texttt{label\_noise\_grokking}]
\small
\begin{verbatim}
"Name": "label_noise_grokking",
"Title": "Grokking Under Noise: The Impact of Systematic and Random Label
Errors on Algorithmic Learning",
"Experiment": "1. Modify the AbstractDataset class to introduce two types
of label noise: random (labels changed randomly) and systematic (specific
labels consistently flipped). Add a 'noise_type' parameter
(random/systematic) and 'noise_level' parameter (low: 5%
high: 20%
training set while keeping the validation set clean. 3. Run experiments for
each noise type and level across all datasets. 4. Track metrics: time to
grokking (99%
loss, and model confidence (mean softmax probability of correct class). 5.
Plot learning curves and model confidence for different noise types and
levels, highlighting the grokking point for each. 6. Analyze how different
types and levels of label noise affect the timing, speed, and extent of
grokking across different operations. 7. Compare attention patterns between
noisy and clean training at different stages to understand how the model
adapts to noise.",
"Interestingness": 9,
"Feasibility": 8,
"Novelty": 9,
"novel": true
\end{verbatim}
\end{tcolorbox}

\begin{tcolorbox}[breakable,colback=blue!5!white, colframe=blue!75!black, title=Idea 14/50 - \texttt{compositional\_grokking}]
\small
\begin{verbatim}
"Name": "compositional_grokking",
"Title": "Compositional Grokking: Investigating the Relationship Between
Grokking and Compositional Learning in Modular Arithmetic",
"Experiment": "1. Modify ModSumDataset and ModSubtractDataset to include
composite operations: (a + b) - c mod p and (a - b) + c mod p. 2. Implement
new dataset class CompositeModDataset for these operations. 3. Run
experiments comparing learning curves for basic (a + b, a - b) and
composite operations. 4. Track metrics: time to grokking for basic vs.
composite operations, correlation between grokking times, final accuracies.
5. Analyze attention patterns to see if the model learns to attend to
intermediate results in composite operations. 6. Implement a 'compositional
generalization' test by training on basic operations and testing on their
compositions. 7. Compare internal representations (e.g., using PCA on
hidden states) for basic vs. composite operations at different stages of
training.",
"Interestingness": 9,
"Feasibility": 6,
"Novelty": 9,
"novel": true
\end{verbatim}
\end{tcolorbox}

\begin{tcolorbox}[breakable,colback=blue!5!white, colframe=blue!75!black, title=Idea 15/50 - \texttt{mutual\_information\_grokking}]
\small
\begin{verbatim}
"Name": "mutual_information_grokking",
"Title": "Information Dynamics in Grokking: Analyzing Mutual Information
Evolution During Algorithmic Learning",
"Experiment": "1. Implement a function to estimate mutual information using
a binning approach for efficiency. 2. Modify the Transformer class to
output hidden states from the final layer. 3. Update the training loop to
calculate and store mutual information between (a) inputs and outputs, and
(b) final hidden states and outputs, at regular intervals. 4. Run
experiments across all datasets, tracking these mutual information metrics
alongside validation accuracy and training loss. 5. Create plots showing
the evolution of both mutual information metrics over training time,
highlighting the grokking point. 6. Analyze how mutual information trends
relate to grokking by testing specific hypotheses: (a) Rapid increase in
hidden state-output mutual information coincides with grokking, (b) Input-
output mutual information stabilizes post-grokking. 7. Compare mutual
information dynamics between different operations and model sizes to
identify common patterns in successful grokking.",
"Interestingness": 9,
"Feasibility": 6,
"Novelty": 9,
"novel": true
\end{verbatim}
\end{tcolorbox}

\begin{tcolorbox}[breakable,colback=blue!5!white, colframe=blue!75!black, title=Idea 16/50 - \texttt{sparse\_subnetworks\_grokking}]
\small
\begin{verbatim}
"Name": "sparse_subnetworks_grokking",
"Title": "Sparse Subnetworks in Grokking: Investigating the Emergence of
Critical Structures During Algorithmic Learning",
"Experiment": "1. Implement a simple magnitude-based pruning function for
the Transformer model. 2. Modify the training loop to perform pruning at
key points: before training, just before grokking (based on validation
accuracy), and after grokking. 3. For each pruning point, create sparse
networks at different sparsity levels (e.g., 50%
these sparse networks from the original initialization for a fixed number
of steps. 5. Track metrics: validation accuracy of sparse networks,
sparsity level, and grokking speed (if it occurs). 6. Plot the performance
of sparse networks at different sparsity levels and pruning points. 7.
Compare the structure and performance of sparse networks found before and
after grokking across different operations.",
"Interestingness": 9,
"Feasibility": 8,
"Novelty": 9,
"novel": true
\end{verbatim}
\end{tcolorbox}

\begin{tcolorbox}[breakable,colback=blue!5!white, colframe=blue!75!black, title=Idea 17/50 - \texttt{positional\_encoding\_grokking}]
\small
\begin{verbatim}
"Name": "positional_encoding_grokking",
"Title": "Inductive Biases in Grokking: The Impact of Positional Encoding
Schemes on Algorithmic Learning",
"Experiment": "1. Modify the Transformer class to support three positional
encoding schemes: sinusoidal (current), learned embeddings, and a simple
binary encoding (e.g., [0,1,0,1,0] for 'a o b = c'). 2. Implement these
encoding schemes, ensuring they work with the existing sequence length. 3.
Run experiments for each encoding scheme across all datasets, tracking:
time to grokking (99%
training loss, and attention entropy. 4. Analyze how different encoding
schemes affect attention patterns and grokking behavior for each operation
type. 5. Compare generalization capabilities on sequences with shuffled
operands (e.g., 'b o a = c'). 6. Correlate encoding scheme performance with
operation complexity to identify potential interactions between input
representation and task structure. 7. Discuss implications for designing
transformers for specific algorithmic tasks based on findings.",
"Interestingness": 9,
"Feasibility": 9,
"Novelty": 9,
"novel": true
\end{verbatim}
\end{tcolorbox}

\begin{tcolorbox}[breakable,colback=blue!5!white, colframe=blue!75!black, title=Idea 18/50 - \texttt{adversarial\_robustness\_grokking}]
\small
\begin{verbatim}
"Name": "adversarial_robustness_grokking",
"Title": "Adversarial Robustness During Grokking: Tracking Vulnerability
Evolution in Algorithmic Learning",
"Experiment": "1. Implement a simple perturbation method: randomly flip 1-2
bits in the input representation for modular arithmetic, and swap 1-2
elements for permutations. 2. Modify the training loop to generate
perturbed inputs and evaluate model performance on them every 500 steps. 3.
Track metrics: normal validation accuracy, accuracy on perturbed inputs,
and 'robustness gap' (difference between normal and perturbed accuracy). 4.
Plot the evolution of robustness to perturbations alongside the grokking
curve. 5. Compare robustness before, during, and after grokking across
different operations. 6. Analyze examples of successful perturbations at
different stages of training. 7. Investigate potential correlations between
the timing of grokking and changes in robustness to perturbations.",
"Interestingness": 9,
"Feasibility": 8,
"Novelty": 9,
"novel": true
\end{verbatim}
\end{tcolorbox}

\begin{tcolorbox}[breakable,colback=blue!5!white, colframe=blue!75!black, title=Idea 19/50 - \texttt{critical\_periods\_grokking}]
\small
\begin{verbatim}
"Name": "critical_periods_grokking",
"Title": "Critical Periods in Grokking: The Impact of Timed Learning Rate
Spikes on Algorithmic Learning",
"Experiment": "1. Modify the training loop to support learning rate spikes
at specific points (25%
to apply these spikes, increasing the learning rate by 10x for 100 steps.
3. Run experiments for each spike timing across all datasets (modular
arithmetic and permutations), including a control group with no spikes. 4.
Track metrics: time to grokking, final validation accuracy, and 'spike
impact' (change in validation accuracy slope in 500 steps post-spike). 5.
Plot learning curves highlighting spike points and their impacts. 6.
Analyze how spike timing affects grokking across different operations,
comparing modular arithmetic tasks with permutations. 7. Compare attention
patterns immediately before and after impactful spikes. 8. Correlate spike
impact with the stage of learning (pre-grokking, during grokking, post-
grokking) to identify potential critical periods, assessing whether these
periods are task-specific or general across operations.",
"Interestingness": 9,
"Feasibility": 9,
"Novelty": 9,
"novel": true
\end{verbatim}
\end{tcolorbox}

\begin{tcolorbox}[breakable,colback=blue!5!white, colframe=blue!75!black, title=Idea 20/50 - \texttt{lottery\_tickets\_grokking}]
\small
\begin{verbatim}
"Name": "lottery_tickets_grokking",
"Title": "Lottery Tickets in Grokking: Investigating Sparse Subnetworks
Capable of Algorithmic Learning",
"Experiment": "1. Implement an iterative magnitude pruning function for the
Transformer model. 2. Modify the training loop to support multiple rounds
of train-prune-reset cycles. 3. For each dataset, run experiments with
pruning levels of 30%
iteration, train the network to convergence, prune the specified percentage
of smallest weights, then reset remaining weights to their initial values.
5. Track metrics for each sparse network: time to grokking (or maximum
training time if grokking doesn't occur), final validation accuracy, and
training loss. 6. Introduce a 'grokking efficiency' metric: the ratio of
time to grokking for the sparse network vs. the dense network. For networks
that don't grok, use the maximum training time. 7. Plot learning curves for
each pruning level, highlighting grokking points and grokking efficiency.
8. Compare the structure of sparse networks that achieve grokking across
different operations, focusing on the distribution of preserved weights in
different layers and attention heads. 9. Analyze the correlation between
pruning level and grokking efficiency across different algorithmic tasks,
including cases where grokking fails to occur.",
"Interestingness": 9,
"Feasibility": 8,
"Novelty": 8,
"novel": false
\end{verbatim}
\end{tcolorbox}

\begin{tcolorbox}[breakable,colback=blue!5!white, colframe=blue!75!black, title=Idea 21/50 - \texttt{algebraic\_structure\_grokking}]
\small
\begin{verbatim}
"Name": "algebraic_structure_grokking",
"Title": "Grokking and Algebraic Structure: How Group Properties Influence
Neural Network Learning",
"Experiment": "1. Implement new dataset classes for modular multiplication
and division (modulo p, where p is prime, ensuring proper group
structures). 2. For each operation (addition, multiplication, division),
calculate and store two properties: group order and number of generators.
3. Run experiments for each operation type, tracking: time to grokking,
final validation accuracy, and the two group properties. 4. Plot learning
curves and grokking points for each operation, labeled with their group
properties. 5. Analyze the correlation between group properties and
grokking behavior (e.g., time to grokking, steepness of accuracy
improvement). 6. Compare attention patterns across operations, focusing on
how they reflect the underlying group structure (e.g., uniformity for
commutative operations). 7. Test the model's ability to generalize by
evaluating on compositions of learned operations (e.g., a * b + c mod p)
after training on individual operations.",
"Interestingness": 9,
"Feasibility": 8,
"Novelty": 9,
"novel": true
\end{verbatim}
\end{tcolorbox}

\begin{tcolorbox}[breakable,colback=blue!5!white, colframe=blue!75!black, title=Idea 22/50 - \texttt{mdl\_grokking}]
\small
\begin{verbatim}
"Name": "mdl_grokking",
"Title": "Minimum Description Length and Grokking: Investigating the
Relationship Between Model Compression and Algorithmic Learning",
"Experiment": "1. Implement functions to calculate model complexity: (a) L2
norm of weights, (b) number of bits to store parameters at different
precisions, (c) effective number of parameters using BIC. 2. Modify the
training loop to track these complexity measures alongside existing
metrics. 3. Run experiments across all datasets, recording complexity
measures, validation accuracy, and training loss at regular intervals. 4.
Plot the evolution of model complexity alongside the grokking curve. 5.
Analyze the correlation between sudden decreases in model complexity and
the onset of grokking, including statistical tests for significance. 6.
Compare complexity dynamics across different operations and model sizes. 7.
Visualize weight distributions at pre-grokking, during grokking, and post-
grokking stages. 8. Implement and compare two early stopping mechanisms:
one based on model complexity stabilization and another based on validation
loss stabilization.",
"Interestingness": 9,
"Feasibility": 8,
"Novelty": 9,
"novel": true
\end{verbatim}
\end{tcolorbox}

\begin{tcolorbox}[breakable,colback=blue!5!white, colframe=blue!75!black, title=Idea 23/50 - \texttt{invariance\_learning\_grokking}]
\small
\begin{verbatim}
"Name": "invariance_learning_grokking",
"Title": "Learning Invariances in Grokking: Tracking Symmetry Awareness
During Algorithmic Learning",
"Experiment": "1. Modify AbstractDataset to generate transformed versions
of inputs (cyclic shifts for modular arithmetic, relabelings for
permutations). 2. Update the evaluation function to test model predictions
on both original and transformed inputs. 3. Implement an 'invariance score'
metric: mean absolute difference between predictions on original and
transformed inputs. 4. Modify the training loop to calculate and store the
invariance score at regular intervals. 5. Run experiments across all
datasets, tracking the invariance score alongside existing metrics. 6. Plot
the evolution of the invariance score alongside the grokking curve. 7.
Analyze how the invariance score changes before, during, and after
grokking. 8. Compare invariance learning across different operations and
model sizes. 9. Investigate correlation between invariance score and
generalization performance.",
"Interestingness": 9,
"Feasibility": 8,
"Novelty": 9,
"novel": true
\end{verbatim}
\end{tcolorbox}

\begin{tcolorbox}[breakable,colback=blue!5!white, colframe=blue!75!black, title=Idea 24/50 - \texttt{grokking\_double\_descent}]
\small
\begin{verbatim}
"Name": "grokking_double_descent",
"Title": "Grokking and Double Descent: Exploring the Intersection of Two
Deep Learning Phenomena",
"Experiment": "1. Create a range of model sizes by varying num_layers (1 to
8) and dim_model (32 to 512). 2. For each dataset, train models of
different sizes, tracking validation accuracy, training loss, and time to
grokking (99%
parameters to identify double descent behavior. 4. On the same plot, mark
the point where grokking occurs for each model size. 5. Analyze the
relationship between grokking timing and the different regimes of the
double descent curve (under-parameterized, critical, over-parameterized).
6. Calculate the correlation between model size and time to grokking. 7.
Compare double descent and grokking behavior across different operations
(modular arithmetic vs. permutations). 8. Investigate whether grokking
consistently occurs in a specific regime of the double descent curve.",
"Interestingness": 9,
"Feasibility": 8,
"Novelty": 9,
"novel": false
\end{verbatim}
\end{tcolorbox}

\begin{tcolorbox}[breakable,colback=blue!5!white, colframe=blue!75!black, title=Idea 25/50 - \texttt{ntk\_alignment\_grokking}]
\small
\begin{verbatim}
"Name": "ntk_alignment_grokking",
"Title": "NTK-Output Alignment in Grokking: Tracking Feature Learning
Dynamics in Algorithmic Tasks",
"Experiment": "1. Implement a function to compute the NTK-output alignment:
the cosine similarity between the NTK's top eigenvector and the output
gradient. 2. Modify the training loop to compute and store this alignment
metric every 100 steps. 3. Run experiments across all datasets, tracking
NTK-output alignment alongside validation accuracy and training loss. 4.
Plot the evolution of NTK-output alignment alongside the grokking curve. 5.
Analyze how the alignment changes before, during, and after grokking,
identifying any consistent patterns across different operations. 6.
Investigate correlations between sudden changes in alignment and the onset
of grokking. 7. Compare alignment dynamics for models that achieve grokking
vs. those that don't. 8. Experiment with using the alignment metric as an
early stopping criterion or to adjust learning rates dynamically. 9.
Discuss implications of findings for understanding feature learning and
generalization in grokking.",
"Interestingness": 9,
"Feasibility": 8,
"Novelty": 9,
"novel": true
\end{verbatim}
\end{tcolorbox}

\begin{tcolorbox}[breakable,colback=blue!5!white, colframe=blue!75!black, title=Idea 26/50 - \texttt{loss\_landscape\_grokking}]
\small
\begin{verbatim}
"Name": "loss_landscape_grokking",
"Title": "Loss Landscape Evolution in Grokking: Geometric Insights into
Algorithmic Learning",
"Experiment": "1. Implement functions to compute and visualize 2D loss
landscapes using filter-wise normalization. 2. Modify the training loop to
save model checkpoints at key points: start of training, 25%
just before grokking (based on validation accuracy), during grokking, and
after grokking. 3. For each checkpoint, compute and store 2D loss landscape
visualizations. 4. Define quantitative metrics for loss landscape
characteristics: (a) local smoothness (average gradient magnitude), (b)
global convexity (ratio of loss at edges to center), (c) barrier height
(maximum loss along minimum loss path). 5. Run experiments across all
datasets, generating loss landscapes and computing metrics at key points.
6. Create side-by-side comparisons of loss landscapes at different stages
of training for each operation. 7. Analyze how loss landscape metrics
change before, during, and after grokking. 8. Compare loss landscape
evolution between operations that grok quickly vs. slowly. 9. Investigate
correlations between changes in loss landscape metrics and the onset of
grokking.",
"Interestingness": 9,
"Feasibility": 8,
"Novelty": 8,
"novel": true
\end{verbatim}
\end{tcolorbox}

\begin{tcolorbox}[breakable,colback=blue!5!white, colframe=blue!75!black, title=Idea 27/50 - \texttt{neural\_collapse\_grokking}]
\small
\begin{verbatim}
"Name": "neural_collapse_grokking",
"Title": "Neural Collapse in Grokking: Investigating Feature Geometry
During Algorithmic Learning",
"Experiment": "1. Modify Transformer to output final layer features. 2.
Implement functions to compute class means and covariances. 3. Calculate
simplified neural collapse metrics: (a) average cosine similarity between
class means, (b) ratio of within-class to between-class variances. 4. Track
these metrics every 500 steps during training. 5. Run experiments on
modular arithmetic and permutation datasets. 6. Plot neural collapse
metrics alongside grokking curves. 7. Analyze changes in metrics before,
during, and after grokking. 8. Compare neural collapse dynamics between
operations that grok quickly vs. slowly. 9. Visualize class mean
trajectories in 2D/3D using PCA. 10. Discuss implications for understanding
both grokking and general neural network learning dynamics.",
"Interestingness": 9,
"Feasibility": 6,
"Novelty": 9,
"novel": true
\end{verbatim}
\end{tcolorbox}

\begin{tcolorbox}[breakable,colback=blue!5!white, colframe=blue!75!black, title=Idea 28/50 - \texttt{data\_augmentation\_grokking}]
\small
\begin{verbatim}
"Name": "data_augmentation_grokking",
"Title": "Data Augmentation in Grokking: The Impact of Input
Transformations on Algorithmic Learning",
"Experiment": "1. Implement task-specific augmentations: (a) For modular
arithmetic: add random offsets (mod p) to inputs. (b) For permutations:
apply random permutations to inputs and outputs. 2. Modify GroupDataset to
apply augmentations with 0%
for each augmentation level across all datasets. 4. Track metrics: time to
grokking (99%
'augmentation generalization gap' (difference between augmented and non-
augmented validation accuracy). 5. Plot learning curves and generalization
gaps for each augmentation level. 6. Analyze the correlation between
augmentation level and grokking speed. 7. Compare attention patterns
between augmentation levels to understand representation changes. 8.
Discuss implications for designing data augmentation strategies in
algorithmic learning tasks.",
"Interestingness": 9,
"Feasibility": 9,
"Novelty": 8,
"novel": true
\end{verbatim}
\end{tcolorbox}

\begin{tcolorbox}[breakable,colback=blue!5!white, colframe=blue!75!black, title=Idea 29/50 - \texttt{emergent\_grokking}]
\small
\begin{verbatim}
"Name": "emergent_grokking",
"Title": "Emergent Abilities in Grokking: Investigating Scale-Dependent
Algorithmic Learning",
"Experiment": "1. Modify existing datasets to include 'simple' and
'complex' versions (e.g., mod sum with small vs. large primes). 2. Adjust
Transformer class to scale from tiny (1 layer, 64 dim) to medium (4 layers,
512 dim). 3. For each operation, train models of increasing size, tracking
grokking time and performance on both simple and complex versions. 4.
Implement a generalization test for each operation (e.g., mod sum with even
larger primes). 5. Plot learning curves for different model sizes,
highlighting grokking points. 6. Create heatmaps of model size vs.
operation complexity, showing grokking time and generalization test
results. 7. Perform statistical analysis to identify significant jumps in
performance across model sizes, using metrics such as accuracy increase
rate and time to reach 99%
behavior patterns across different operation types.",
"Interestingness": 9,
"Feasibility": 8,
"Novelty": 9,
"novel": true
\end{verbatim}
\end{tcolorbox}

\begin{tcolorbox}[breakable,colback=blue!5!white, colframe=blue!75!black, title=Idea 30/50 - \texttt{functional\_modularity\_grokking}]
\small
\begin{verbatim}
"Name": "functional_modularity_grokking",
"Title": "Functional Modularity in Grokking: Analyzing Emergent
Specialization in Transformer Networks During Algorithmic Learning",
"Experiment": "1. Implement functions to track weight update patterns and
attention focus for each layer and head. 2. Modify the training loop to
compute and store these metrics at regular intervals. 3. Define a
'functional modularity score' based on the consistency of weight updates
and attention patterns for specific input types. 4. Run experiments across
all datasets, tracking the functional modularity score alongside existing
metrics. 5. Plot the evolution of functional modularity alongside the
grokking curve. 6. Analyze how functional modularity changes before,
during, and after grokking. 7. Visualize the most consistent patterns at
different stages of training and interpret their functions. 8. Compare
functional modularity dynamics between different operations and model
sizes. 9. Investigate correlations between functional modularity and
grokking speed or generalization performance.",
"Interestingness": 9,
"Feasibility": 8,
"Novelty": 9,
"novel": true
\end{verbatim}
\end{tcolorbox}

\begin{tcolorbox}[breakable,colback=blue!5!white, colframe=blue!75!black, title=Idea 31/50 - \texttt{information\_compression\_grokking}]
\small
\begin{verbatim}
"Name": "information_compression_grokking",
"Title": "Information Compression in Grokking: Analyzing Representational
Dynamics During Algorithmic Learning",
"Experiment": "1. Modify Transformer class to include a bottleneck layer
(smaller dimension linear layer) after the encoder. 2. Implement function
to compute activation sparsity (%
bottleneck layer. 3. Update training loop to compute and store activation
sparsity and gradient magnitudes of the bottleneck layer at regular
intervals. 4. Run experiments with different bottleneck sizes (e.g., 25%
50%
to grokking, final validation accuracy, activation sparsity, and gradient
magnitudes. 6. Plot activation sparsity and gradient magnitude evolution
alongside grokking curves for each bottleneck size. 7. Analyze how these
metrics change before, during, and after grokking. 8. Test generalization
by evaluating models on slightly out-of-distribution examples (e.g., larger
numbers in modular arithmetic). 9. Investigate correlation between optimal
compression (measured by activation sparsity) and grokking speed,
generalization performance.",
"Interestingness": 9,
"Feasibility": 8,
"Novelty": 8,
"novel": true
\end{verbatim}
\end{tcolorbox}

\begin{tcolorbox}[breakable,colback=blue!5!white, colframe=blue!75!black, title=Idea 32/50 - \texttt{critical\_learning\_periods\_grokking}]
\small
\begin{verbatim}
"Name": "critical_learning_periods_grokking",
"Title": "Critical Learning Periods in Grokking: Temporal Dynamics of
Algorithmic Understanding",
"Experiment": "1. Modify the training loop to support 'intervention
periods' where learning rate is increased by 5x for 100 steps. 2. Implement
a sliding window intervention strategy, with windows of 500 steps, starting
every 250 steps. 3. Run experiments for each window across all datasets and
three model sizes (small, medium, large), including a control group with no
interventions. 4. Track metrics: time to grokking, final validation
accuracy, and 'intervention impact' (area under the validation accuracy
curve for 500 steps post-intervention). 5. Plot learning curves
highlighting intervention windows and their impacts. 6. Create heatmaps
visualizing intervention impact across time windows and model sizes for
each operation. 7. Analyze how intervention timing affects grokking across
different operations and model sizes. 8. Compare attention patterns
immediately before and after impactful interventions. 9. Investigate
whether certain operations or model sizes have more pronounced critical
periods than others. 10. Discuss implications for curriculum design in
machine learning and potential applications in continual and transfer
learning.",
"Interestingness": 9,
"Feasibility": 7,
"Novelty": 9,
"novel": true
\end{verbatim}
\end{tcolorbox}

\begin{tcolorbox}[breakable,colback=blue!5!white, colframe=blue!75!black, title=Idea 33/50 - \texttt{simplicity\_bias\_grokking}]
\small
\begin{verbatim}
"Name": "simplicity_bias_grokking",
"Title": "Simplicity Bias in Grokking: Analyzing Weight Matrix Complexity
During Algorithmic Learning",
"Experiment": "1. Modify AbstractDataset to include two complexity levels
for each operation (e.g., small vs. large prime for modular arithmetic,
short vs. long permutations). 2. Implement a function to compute the
effective rank of weight matrices using singular value decomposition. 3.
Update the training loop to compute and store the effective rank for each
layer every 500 steps. 4. Run experiments across all datasets and both
complexity levels, tracking effective rank alongside existing metrics. 5.
Plot the evolution of effective rank alongside grokking curves for each
complexity level and operation. 6. Analyze how effective rank changes
before, during, and after grokking, and how this relates to task
complexity. 7. Investigate correlations between effective rank dynamics and
grokking speed or generalization performance. 8. Compare effective rank
patterns across different operations and model sizes. 9. Contrast effective
rank dynamics between operations that grok quickly versus those that grok
slowly or fail to grok. 10. Experiment with using effective rank as an
indicator for the onset of grokking.",
"Interestingness": 9,
"Feasibility": 8,
"Novelty": 9,
"novel": true
\end{verbatim}
\end{tcolorbox}

\begin{tcolorbox}[breakable,colback=blue!5!white, colframe=blue!75!black, title=Idea 34/50 - \texttt{lucky\_initializations\_grokking}]
\small
\begin{verbatim}
"Name": "lucky_initializations_grokking",
"Title": "Lucky Initializations in Grokking: Identifying and Analyzing
Favorable Starting Points for Algorithmic Learning",
"Experiment": "1. Implement a function to generate and store 50 random
initializations for the Transformer model. 2. Modify the training loop to
support training from stored initializations and different learning rates.
3. For each dataset, train models from the 50 initializations with 3
learning rates, tracking 'grokking efficiency' (ratio of validation
accuracy to training steps at 99%
initializations (top 20%
characteristics of lucky initializations: weight distribution statistics,
layerwise norms, and attention pattern initialization. 6. Implement a
function to visualize the loss landscape around initial points using
filter-wise normalization. 7. Compare lucky initializations across
different operations to identify common patterns. 8. Develop a simple
predictor for initialization 'luckiness' based on identified
characteristics. 9. Test transfer of lucky initializations across tasks and
learning rates.",
"Interestingness": 9,
"Feasibility": 9,
"Novelty": 9,
"novel": true
\end{verbatim}
\end{tcolorbox}

\begin{tcolorbox}[breakable,colback=blue!5!white, colframe=blue!75!black, title=Idea 35/50 - \texttt{relative\_attention\_grokking}]
\small
\begin{verbatim}
"Name": "relative_attention_grokking",
"Title": "Relative Positional Attention and Its Impact on Grokking in
Algorithmic Learning",
"Experiment": "1. Modify the DecoderBlock class to support two attention
types: standard (current) and relative positional. 2. Implement relative
positional attention, ensuring it works with the existing sequence length.
3. Update the Transformer class to accept an attention_type parameter. 4.
Run experiments for both attention types across all datasets, tracking:
time to grokking (99%
training loss, grokking transition sharpness (rate of validation accuracy
increase), and post-grokking stability (variance in validation accuracy
after reaching 99%
highlighting grokking points and transition periods. 6. Visualize and
compare attention patterns between the two mechanisms at key stages: pre-
grokking, during grokking transition, and post-grokking. 7. Analyze how
relative positional attention affects grokking behavior, transition
sharpness, and stability for each operation type compared to standard
attention. 8. Investigate correlations between attention type and grokking
speed or post-grokking stability. 9. Discuss implications for designing
transformers for specific algorithmic tasks based on findings.",
"Interestingness": 9,
"Feasibility": 8,
"Novelty": 8,
"novel": true
\end{verbatim}
\end{tcolorbox}

\begin{tcolorbox}[breakable,colback=blue!5!white, colframe=blue!75!black, title=Idea 36/50 - \texttt{grokking\_task\_interference}]
\small
\begin{verbatim}
"Name": "grokking_task_interference",
"Title": "Grokking and Task Interference: Exploring the Stability of
Algorithmic Understanding",
"Experiment": "1. Modify the training loop to support learning two modular
arithmetic operations sequentially (e.g., addition then multiplication). 2.
Implement a task scheduler that switches between tasks at regular
intervals. 3. Create a 'dual-task evaluation' function to assess
performance on both tasks simultaneously. 4. Track metrics: time to
grokking for each task, performance on the first task while learning the
second, and a 'grokking stability' score (maintenance of >95%
task 1 while learning task 2). 5. Run experiments with different task
switching frequencies. 6. Analyze how grokking on one task affects learning
speed and grokking on the subsequent task. 7. Visualize attention patterns
before and after introducing the second task to understand representation
changes. 8. Investigate the correlation between grokking speed on the first
task and stability of that understanding when learning the second task. 9.
Compare results with a baseline of learning both tasks simultaneously to
isolate the effects of sequential learning.",
"Interestingness": 9,
"Feasibility": 8,
"Novelty": 9,
"novel": true
\end{verbatim}
\end{tcolorbox}

\begin{tcolorbox}[breakable,colback=blue!5!white, colframe=blue!75!black, title=Idea 37/50 - \texttt{attention\_inductive\_bias\_grokking}]
\small
\begin{verbatim}
"Name": "attention_inductive_bias_grokking",
"Title": "Inductive Biases in Attention Mechanisms: Their Impact on
Grokking in Algorithmic Learning",
"Experiment": "1. Modify DecoderBlock class to support two attention
mechanisms: standard dot-product and additive (Bahdanau). 2. Implement
these attention mechanisms, ensuring compatibility with existing
architecture. 3. Update Transformer class to accept an attention_type
parameter. 4. Select a subset of most illustrative datasets based on
preliminary experiments. 5. Run experiments for each attention type on
selected datasets, tracking: time to grokking (99%
final validation accuracy, training loss, and 'grokking transition
sharpness' (defined as the maximum rate of validation accuracy increase
over any 500-step window). 6. Implement a simple generalization test using
slightly out-of-distribution examples (e.g., larger numbers for modular
arithmetic). 7. Plot learning curves for each attention type, highlighting
grokking points and transition periods. 8. Analyze how different attention
mechanisms affect grokking behavior, transition sharpness, and
generalization performance for each operation type. 9. Visualize attention
patterns for each mechanism at key stages: pre-grokking, during grokking
transition, and post-grokking. 10. Discuss implications for designing
transformers with appropriate inductive biases for specific types of
algorithmic learning tasks.",
"Interestingness": 9,
"Feasibility": 9,
"Novelty": 9,
"novel": true
\end{verbatim}
\end{tcolorbox}

\begin{tcolorbox}[breakable,colback=blue!5!white, colframe=blue!75!black, title=Idea 38/50 - \texttt{gradient\_dynamics\_grokking}]
\small
\begin{verbatim}
"Name": "gradient_dynamics_grokking",
"Title": "Gradient Dynamics in Grokking: Analyzing Information Flow
Efficiency During Algorithmic Learning",
"Experiment": "1. Modify the training loop to compute gradient statistics
(sparsity and magnitude distribution) for each layer. 2. Implement
functions to calculate gradient sparsity (%
magnitude percentiles. 3. Update training process to store these metrics
every 500 steps. 4. Run experiments across all datasets, tracking gradient
metrics alongside existing performance metrics. 5. Plot the evolution of
gradient sparsity and magnitude distributions alongside grokking curves. 6.
Analyze how gradient dynamics change before, during, and after grokking. 7.
Compare gradient patterns between operations that grok quickly vs. slowly.
8. Investigate correlations between changes in gradient dynamics and
grokking speed or generalization performance. 9. Visualize gradient flow
patterns at key stages: pre-grokking, during grokking transition, and post-
grokking.",
"Interestingness": 9,
"Feasibility": 8,
"Novelty": 8,
"novel": true
\end{verbatim}
\end{tcolorbox}

\begin{tcolorbox}[breakable,colback=blue!5!white, colframe=blue!75!black, title=Idea 39/50 - \texttt{adaptive\_curriculum\_grokking}]
\small
\begin{verbatim}
"Name": "adaptive_curriculum_grokking",
"Title": "Adaptive Curriculum Learning in Grokking: Optimizing Example
Difficulty for Efficient Algorithmic Understanding",
"Experiment": "1. Modify AbstractDataset to include a difficulty scoring
function (e.g., input magnitude for modular arithmetic, cycle length for
permutations). 2. Implement adaptive sampling strategy: start with easiest
20%
level exceeds 90%
tracking difficulty of selected examples. 4. Run experiments comparing
adaptive curriculum, random sampling, and static curriculum (increasing
difficulty linearly) across all datasets. 5. Track metrics: time to
grokking, final validation accuracy, learning trajectory smoothness, and
example difficulty distribution over time. 6. Analyze relationship between
difficulty progression and grokking onset. 7. Visualize learning curves and
difficulty progression for each strategy. 8. Compare consistency and speed
of grokking across different random seeds for each strategy. 9. Analyze
computational efficiency by comparing total number of examples needed to
achieve grokking for each strategy. 10. Compare attention patterns at key
points (pre-grokking, during grokking, post-grokking) across strategies to
understand how adaptive curriculum affects internal representations.",
"Interestingness": 9,
"Feasibility": 8,
"Novelty": 9,
"novel": true
\end{verbatim}
\end{tcolorbox}

\begin{tcolorbox}[breakable,colback=blue!5!white, colframe=blue!75!black, title=Idea 40/50 - \texttt{task\_structure\_grokking}]
\small
\begin{verbatim}
"Name": "task_structure_grokking",
"Title": "Task Structure and Grokking: Investigating the Relationship
Between Algorithmic Complexity and Learning Dynamics",
"Experiment": "1. Modify AbstractDataset to include a
'structural_complexity' score based on: a) number of unique outputs, b)
input-output correlation, c) algebraic degree for modular operations or
cycle structure for permutations. 2. Extend existing dataset classes to
include a wider range of operations (e.g., modular addition,
multiplication, exponentiation; simple and complex permutations). 3. Run
experiments across all operations, tracking time to grokking, final
validation accuracy, and learning curve smoothness. 4. Plot grokking
metrics against structural complexity scores, comparing trends between
modular arithmetic and permutation tasks. 5. Analyze correlation between
structural complexity and grokking behavior. 6. Compare attention patterns
and gradient flows across tasks of different complexity. 7. Implement a
generalization test where models trained on simpler structures are
evaluated on more complex ones. 8. Discuss implications for neural network
learning on structured vs. unstructured tasks in general machine learning
contexts.",
"Interestingness": 9,
"Feasibility": 9,
"Novelty": 9,
"novel": true
\end{verbatim}
\end{tcolorbox}

\begin{tcolorbox}[breakable,colback=blue!5!white, colframe=blue!75!black, title=Idea 41/50 - \texttt{numerical\_base\_grokking}]
\small
\begin{verbatim}
"Name": "numerical_base_grokking",
"Title": "Numerical Base and Grokking: How Input Representation Affects
Pattern Recognition in Algorithmic Learning",
"Experiment": "1. Modify AbstractDataset and modular arithmetic dataset
classes to support binary and decimal bases. 2. Implement functions to
convert between bases and adjust the encode/decode methods. 3. Update the
Transformer class to handle variable input lengths. 4. Run experiments for
binary and decimal bases on modular addition and multiplication tasks. 5.
Track metrics: time to grokking (99%
accuracy, training loss, and 'cross-base generalization' (accuracy when
testing on the other base). 6. Plot learning curves for each base,
highlighting grokking points. 7. Compare learning curves for binary (0-3)
vs decimal (0-9) to isolate base effects from sequence length. 8. Analyze
how different bases affect grokking speed and pattern recognition. 9.
Compare attention patterns across bases at key stages: pre-grokking, during
grokking, and post-grokking. 10. Discuss implications for choosing input
representations in mathematical machine learning tasks.",
"Interestingness": 9,
"Feasibility": 9,
"Novelty": 9,
"novel": true
\end{verbatim}
\end{tcolorbox}

\begin{tcolorbox}[breakable,colback=blue!5!white, colframe=blue!75!black, title=Idea 42/50 - \texttt{activation\_function\_grokking}]
\small
\begin{verbatim}
"Name": "activation_function_grokking",
"Title": "Activation Functions and Grokking: Investigating the Role of Non-
linearity in Algorithmic Learning and Generalization",
"Experiment": "1. Modify the DecoderBlock class to support multiple
activation functions (ReLU, GELU, Tanh). 2. Update the Transformer class to
accept an activation_type parameter, allowing for both uniform and hybrid
activation setups. 3. Run experiments comparing the baseline (GELU) with
ReLU, Tanh, and a hybrid setup (ReLU in lower layers, Tanh in upper layers)
across all datasets. 4. Track metrics: time to grokking (99%
accuracy), final validation accuracy, training loss, 'grokking transition
sharpness', and gradient flow statistics. 5. Plot learning curves for each
activation setup, highlighting grokking points and transition periods. 6.
Visualize decision boundaries at different training stages for each
activation setup. 7. Analyze how different activation functions affect
grokking behavior, transition sharpness, and final performance for each
operation type. 8. Compare hidden representations (using t-SNE) across
activation setups at key stages: pre-grokking, during grokking transition,
and post-grokking. 9. Investigate the relationship between activation
function properties and the trade-off between memorization and
generalization. 10. Discuss implications for choosing activation functions
in tasks requiring pattern discovery and generalization beyond algorithmic
learning.",
"Interestingness": 9,
"Feasibility": 9,
"Novelty": 9,
"novel": true
\end{verbatim}
\end{tcolorbox}

\begin{tcolorbox}[breakable,colback=blue!5!white, colframe=blue!75!black, title=Idea 43/50 - \texttt{phase\_transition\_grokking}]
\small
\begin{verbatim}
"Name": "phase_transition_grokking",
"Title": "Grokking as a Phase Transition: Characterizing Critical Behavior
in Algorithmic Learning",
"Experiment": "1. Implement functions to track key metrics: validation
accuracy, training loss, gradient norm, and weight norm. 2. Modify training
loop to compute and store these metrics every 100 steps. 3. Run experiments
across all datasets, with finer-grained tracking (every 10 steps) around
the suspected grokking point. 4. Implement analysis tools to detect sudden
changes or discontinuities in metrics. 5. Plot all metrics on a single,
multi-axis graph to visualize potential phase transitions. 6. Calculate
susceptibility using fluctuations in validation accuracy near the grokking
point. 7. Analyze scaling behavior of susceptibility to identify critical
exponents, if any. 8. Compare phase transition characteristics across
different operations and model sizes. 9. Investigate whether manipulating
learning rate or gradient clipping can induce or prevent grokking phase
transitions.",
"Interestingness": 9,
"Feasibility": 8,
"Novelty": 9,
"novel": false
\end{verbatim}
\end{tcolorbox}

\begin{tcolorbox}[breakable,colback=blue!5!white, colframe=blue!75!black, title=Idea 44/50 - \texttt{effective\_dimension\_grokking}]
\small
\begin{verbatim}
"Name": "effective_dimension_grokking",
"Title": "Effective Dimension Dynamics in Grokking: Analyzing
Representational Complexity During Algorithmic Learning",
"Experiment": "1. Implement functions to compute the rank and top-k
singular values of weight matrices. 2. Modify the training loop to compute
and store these metrics every 500 steps for each layer. 3. Run experiments
across all datasets, tracking rank and singular value distributions
alongside existing performance metrics. 4. Implement a simple MLP baseline
that doesn't exhibit grokking for comparison. 5. Plot the evolution of rank
and singular value distributions alongside grokking curves for both
Transformer and MLP models. 6. Analyze how these metrics change before,
during, and after grokking in the Transformer, contrasting with the MLP. 7.
Compare rank dynamics between operations that grok quickly vs. slowly. 8.
Investigate correlations between changes in rank/singular values and
grokking speed or generalization performance. 9. Visualize the relationship
between these metrics and other performance indicators at different stages
of training.",
"Interestingness": 9,
"Feasibility": 8,
"Novelty": 9,
"novel": true
\end{verbatim}
\end{tcolorbox}

\begin{tcolorbox}[breakable,colback=blue!5!white, colframe=blue!75!black, title=Idea 45/50 - \texttt{representation\_entropy\_grokking}]
\small
\begin{verbatim}
"Name": "representation_entropy_grokking",
"Title": "Representation Entropy in Grokking: Tracking the Simplification
of Learned Concepts",
"Experiment": "1. Implement a function to compute the entropy of the
model's internal representations. 2. Modify the Transformer class to output
intermediate representations. 3. Update the training loop to compute and
store the representation entropy every 500 steps. 4. Run experiments across
all datasets, including configurations that lead to successful grokking and
those that don't (e.g., by varying model size or learning rate). 5. Track
entropy alongside existing performance metrics. 6. Plot the evolution of
representation entropy alongside grokking curves for both successful and
unsuccessful cases. 7. Analyze how representation entropy changes before,
during, and after grokking in successful cases, and compare with
unsuccessful cases. 8. Investigate correlations between changes in
representation entropy and grokking speed or generalization performance. 9.
Visualize the relationship between entropy and other performance indicators
at different stages of training. 10. Plot entropy distributions across
different layers of the model to understand how different parts contribute
to concept simplification.",
"Interestingness": 9,
"Feasibility": 8,
"Novelty": 8,
"novel": true
\end{verbatim}
\end{tcolorbox}

\begin{tcolorbox}[breakable,colback=blue!5!white, colframe=blue!75!black, title=Idea 46/50 - \texttt{mutual\_information\_grokking}]
\small
\begin{verbatim}
"Name": "mutual_information_grokking",
"Title": "Mutual Information Dynamics in Grokking: Tracing Information Flow
During Algorithmic Learning",
"Experiment": "1. Modify Transformer class to output representations from
input embedding, middle layer, and final layer. 2. Implement MINE (Mutual
Information Neural Estimation) for efficient mutual information
approximation. 3. Update training loop to compute and store mutual
information estimates between input-middle, input-output, and middle-output
every 500 steps. 4. Run experiments across all datasets, tracking mutual
information alongside existing performance metrics. 5. Plot the evolution
of mutual information alongside grokking curves and generalization gap. 6.
Analyze how mutual information changes before, during, and after grokking,
particularly in relation to the generalization gap. 7. Compare mutual
information dynamics between operations that grok quickly vs. slowly. 8.
Investigate correlations between changes in mutual information and grokking
speed or generalization performance.",
"Interestingness": 9,
"Feasibility": 8,
"Novelty": 8,
"novel": true
\end{verbatim}
\end{tcolorbox}

\begin{tcolorbox}[breakable,colback=blue!5!white, colframe=blue!75!black, title=Idea 47/50 - \texttt{lottery\_tickets\_grokking}]
\small
\begin{verbatim}
"Name": "lottery_tickets_grokking",
"Title": "Lottery Tickets in Grokking: Sparse Subnetworks and Sudden
Generalization",
"Experiment": "1. Implement iterative magnitude pruning for the Transformer
model. 2. Modify training loop for train-prune-reset cycles. 3. For each
dataset, run experiments with pruning levels of 50%
iterations. 4. Track metrics: time to grokking, final validation accuracy,
training loss, and 'grokking efficiency' (ratio of time to grokking for
sparse vs. dense network). 5. Plot learning curves for each pruning level,
highlighting grokking points. 6. Compare sparse network structures that
achieve grokking across operations. 7. Analyze correlation between pruning
level and grokking efficiency. 8. Implement simple MLP baseline without
grokking for comparison. 9. Visualize weight distributions of winning
tickets pre- and post-grokking.",
"Interestingness": 9,
"Feasibility": 9,
"Novelty": 8,
"novel": false
\end{verbatim}
\end{tcolorbox}

\begin{tcolorbox}[breakable,colback=blue!5!white, colframe=blue!75!black, title=Idea 48/50 - \texttt{architecture\_inductive\_bias\_grokking}]
\small
\begin{verbatim}
"Name": "architecture_inductive_bias_grokking",
"Title": "Architectural Inductive Biases and Grokking: Comparing Sudden
Generalization Across Neural Network Types",
"Experiment": "1. Implement simplified 1D CNN and LSTM model classes
compatible with existing sequence-based datasets. 2. Modify training loop
to support multiple model types. 3. Run experiments comparing Transformer,
1D CNN, and LSTM models across modular arithmetic datasets. 4. Track
metrics: time to grokking, final validation accuracy, training loss, and
architecture-specific indicators (attention patterns for Transformer,
filter activations for CNN, forget gate activations for LSTM). 5. Plot
learning curves for each architecture, highlighting grokking points. 6.
Analyze how different architectures affect grokking behavior, speed, and
final performance for each operation type. 7. Compare internal
representations (using t-SNE) across architectures at key stages: pre-
grokking, during grokking transition, and post-grokking. 8. Investigate the
relationship between architectural inductive biases and the trade-off
between memorization and generalization in modular arithmetic tasks.",
"Interestingness": 9,
"Feasibility": 8,
"Novelty": 8,
"novel": true
\end{verbatim}
\end{tcolorbox}

\begin{tcolorbox}[breakable,colback=blue!5!white, colframe=blue!75!black, title=Idea 49/50 - \texttt{shortcut\_learning\_grokking}]
\small
\begin{verbatim}
"Name": "shortcut_learning_grokking",
"Title": "Shortcut Learning and Grokking: The Interplay Between Surface
Patterns and Deep Understanding in Algorithmic Learning",
"Experiment": "1. Modify AbstractDataset to include operation-specific
shortcuts: for modular arithmetic, make the result always even if the first
operand is even; for permutations, always swap the first two elements. 2.
Implement a function to gradually remove these shortcuts over training by
reducing their frequency. 3. Update the training loop to apply the shortcut
removal function. 4. Add a 'shortcut reliance' metric: the accuracy
difference between shortcut-following and shortcut-violating examples. 5.
Run experiments with varying shortcut removal rates across datasets. 6.
Track metrics: time to grokking, final validation accuracy, shortcut
reliance over time, and performance on a shortcut-free test set. 7. Plot
learning curves and shortcut reliance alongside grokking curves. 8. Analyze
how shortcut presence and removal affect grokking timing and quality. 9.
Compare attention patterns between models trained with and without
shortcuts at key stages.",
"Interestingness": 9,
"Feasibility": 9,
"Novelty": 9,
"novel": true
\end{verbatim}
\end{tcolorbox}

\begin{tcolorbox}[breakable,colback=blue!5!white, colframe=blue!75!black, title=Idea 50/50 - \texttt{grokking\_forgetting\_complexity}]
\small
\begin{verbatim}
"Name": "grokking_forgetting_complexity",
"Title": "Grokking and Forgetting: The Interplay of Task Complexity and
Sudden Generalization in Algorithmic Learning",
"Experiment": "1. Modify ModSumDataset to support multiple complexity
levels (e.g., modular addition with increasing prime moduli). 2. Update the
training loop to gradually introduce higher complexity levels while
continuously evaluating on all levels. 3. Implement a 'multi-complexity
evaluation' function to assess performance across all complexity levels
simultaneously. 4. Track metrics: time to grokking for each complexity
level, performance on lower complexity levels when grokking occurs on a
higher level, and a 'complexity forgetting score' (decrease in accuracy on
lower complexity levels). 5. Analyze the correlation between grokking
events and performance changes on other complexity levels. 6. Compare
internal representations (using cosine similarity of hidden states) across
complexity levels before and after grokking events. 7. Investigate trends
in grokking speed across increasing complexity levels. 8. Plot learning
curves for all complexity levels simultaneously, highlighting grokking
points and potential forgetting events. 9. Visualize the evolution of
representation similarities over time using heatmaps.",
"Interestingness": 9,
"Feasibility": 9,
"Novelty": 9,
"novel": true
\end{verbatim}
\end{tcolorbox}

